% Cross-cutting index frames

\footnotesize

\begin{frame}{索引：Case 到待做实验}
\scriptsize
\textbf{用途：}本页把当前样例证据映射到后续正式评估；这里不新增结论，只说明每类现象下一步如何量化。

\begin{tabularx}{\linewidth}{l>{\raggedright\arraybackslash}p{3.2cm}>{\raggedright\arraybackslash}p{3.6cm}>{\raggedright\arraybackslash}X}
\toprule
Case & 当前样例读法 & 无训练线 & 训练相关线 \\
\midrule
A / F & evaluation hit 弱到中等，强依赖窗口 & \code{nt_eval_awareness}: awareness grader、prompt echo、realism edit、sampling-window & \code{tr_training_curves}: 观察 awareness 能力是否随 checkpoint 变化 \\
B / I & policy/refusal 约束可读，但多为 prompt-induced & \code{nt_toy_auditing}: toy hidden objective workflow & \code{tr_hidden_objective_auditing}: trained target organism 上验证 root cause \\
C & language signal 强，但文化 cue 和输出后 token 会混淆 & \code{nt_known_truth_language}, \code{nt_nla_qa}: lead-time 与语言标签恢复 & -- \\
D & 数学/验证/冲突主题可读，不是答案 oracle & \code{nt_known_truth_answer}, \code{nt_nla_qa}: 候选答案恢复 & \code{tr_direct_activation_oracle}: activation-to-answer 上限 \\
E & claim candidates 多，但真伪未审 & \code{nt_confabulation_claim}: truth/support/recurrence + deletion ablation & \code{tr_sae_agreement}: 与 SAE feature evidence 交叉验证 \\
G / H & 任务类型、格式、persona 是强混淆项 & \code{nt_known_truth_eval}: task/style control 与 shuffled baseline & \code{tr_training_curves}: 检查训练是否放大模板化解释 \\
\bottomrule
\end{tabularx}
\end{frame}


\begin{frame}{索引：无训练线优先级}
\scriptsize
\textbf{这些实验可直接使用当前 released Qwen L20 AV/AR。}

\begin{tabularx}{\linewidth}{l>{\raggedright\arraybackslash}p{3.0cm}>{\raggedright\arraybackslash}X}
\toprule
实验 & 主要承接 Case & 必须回答的问题 \\
\midrule
\code{nt_eval_awareness} & A / F & evaluation/review hit 是否高于 prompt echo、format-only 和 deployment-like controls？ \\
\code{nt_confabulation_claim} & E & supported/false/recurring claims 删除后，对 AR cos/MSE 的影响是否不同？ \\
\code{nt_known_truth_eval} & C / D / G / H & NLA explanation 能否恢复已知真值，而不只是读出任务格式？ \\
\code{nt_steganography_quality} & E / D & paraphrase、compression、summary 是否导致异常 reconstruction drop？ \\
\code{nt_toy_auditing} & B / I & agent 使用 NLA 是否更容易定位 toy hidden objective 的 root cause？ \\
\bottomrule
\end{tabularx}

\vspace{0.35em}
\alert{口径：}无训练线只能评价 final released NLA 的诊断价值，不能声称能力来自训练过程，也不能复现 trained hidden-objective benchmark。
\end{frame}


\begin{frame}{索引：训练相关线优先级}
\scriptsize
\textbf{这些实验需要 checkpoint、SAE、oracle 或 trained target organism。}

\begin{tabularx}{\linewidth}{l>{\raggedright\arraybackslash}p{3.3cm}>{\raggedright\arraybackslash}X}
\toprule
实验 & 主要承接 Case & 决定性产物 \\
\midrule
\code{tr_training_curves} & A / G / H & 同一批 activations 上的 FVE、known-truth performance、writing quality 随 checkpoint 变化 \\
\code{tr_sae_agreement} & E / G & NLA claim 与 SAE top features 的 thematic / specific-claim agreement \\
\code{tr_direct_activation_oracle} & D & direct activation-to-answer 是否超过 NLA + grader，或以更低成本持平 \\
\code{tr_hidden_objective_auditing} & B / I & trained hidden-objective target 上，NLA 是否提高 root-cause identification win rate \\
\bottomrule
\end{tabularx}

\vspace{0.35em}
\alert{口径：}只有训练相关线完成后，才可以讨论原文最强主张：训练产生的 NLA 能否支持 automated auditing。
\end{frame}


\begin{frame}{索引：四类容易混淆的信号}
\scriptsize
\textbf{后续实验设计必须把这些机制拆开。}

\begin{tabularx}{\linewidth}{l>{\raggedright\arraybackslash}p{3.1cm}>{\raggedright\arraybackslash}X}
\toprule
机制 & 代表 Case & 不应怎样误读 \\
\midrule
格式/体裁 & G / H / A2 / D4 & 高 cos 或低 cos 不等于 hidden state；可能只是输出格式更易或更难 verbalize \\
语言/文化 cue & C & Chinese hit 不等于提前 language switch；Hangzhou/Taipei 也会触发语言邻域 \\
policy/refusal & B / I & prompt-induced refusal 不等于 trained hidden objective；需要 trained target organism \\
claim specificity & E & AR reconstruction 好不代表 claim 为真；必须做人审和 deletion ablation \\
\bottomrule
\end{tabularx}

\vspace{0.35em}
\alert{使用方式：}读任何 case 时，先定位它属于哪类机制，再看对应 \code{nt_*} 或 \code{tr_*} 实验是否已完成。
\end{frame}
